% !TeX root = ../relatorio-projeto.tex
\chapter{Tema e contexto do projeto}

\section{Orientações de uso}

Este protótipo apresenta um planejamento preenchido com situação hipotética.
Substitua o tema, as premissas e as entregas pelas decisões do seu projeto,
com indicação das fontes e evidências realmente obtidas. Não reutilize o
diagnóstico ilustrativo como descrição de uma instituição real.

O roteiro relaciona a modalidade de produto técnico ou tecnológico ao texto
acadêmico de acompanhamento descrito no manual do projeto
\citeind{manualtccundf}. Seu uso deve ser discutido com o orientador e confirmado
conforme o curso. A apresentação usa os recursos existentes do UnDF-TeX, sem
declarar implementação integral de uma norma específica de relatório.

\section{Tema e delimitação}

Propõe-se planejar um protótipo de sistema para registrar e acompanhar
solicitações de manutenção de equipamentos de um laboratório didático.
O recorte inclui abertura, classificação, atribuição de responsável e
acompanhamento da situação da solicitação. Compra de equipamentos, gestão
patrimonial e integração com sistemas institucionais ficam fora do escopo.

\section{Contexto, público e demanda}

Adota-se, somente como hipótese de trabalho, um laboratório em que solicitações
chegam por canais distintos. Nesse cenário, estudantes comunicariam problemas
e uma equipe técnica organizaria o atendimento. A dispersão dos registros
poderia dificultar o acompanhamento, mas essa hipótese ainda precisa de
diagnóstico: não foram realizados entrevistas, observações ou levantamento
de registros administrativos.

A demanda deverá ser confirmada antes do desenvolvimento. O diagnóstico
previsto identificará os canais existentes, os responsáveis, os dados
necessários e as restrições de acesso. Suas evidências deverão fundamentar
os requisitos, sem incluir dados pessoais além dos necessários.

\chapter{Proposta do projeto}

\section{Objetivo e justificativa}

O objetivo geral é desenvolver e avaliar, em ambiente de demonstração, um
protótipo para acompanhar solicitações de manutenção. Os objetivos específicos
são caracterizar o fluxo, definir requisitos, construir as funções prioritárias
e verificar sua adequação a cenários previamente definidos.

A justificativa a investigar é a possibilidade de centralizar informações
necessárias ao acompanhamento. Benefícios como redução de prazo ou aumento
de satisfação são hipóteses; não constituem resultados deste planejamento.

\section{Solução, requisitos e entregáveis}

A solução prevista é uma interface com cadastro de solicitação, consulta de
situação e atualização pelo responsável. Os requisitos iniciais deverão
contemplar identificação do equipamento, descrição do problema, situação e
histórico de alterações. O acesso será definido por função. Na demonstração,
serão usados equipamentos e usuários fictícios.

As entregas propostas são um diagnóstico documentado, uma lista priorizada
de requisitos, um protótipo executável, registros dos testes e um relatório
final de análise. Nenhuma dessas entregas é declarada concluída neste documento.

\section{Critérios de aceitação previstos}

Para cada requisito priorizado, será definido um cenário com entrada, resultado
esperado e evidência de execução. O protótipo deverá permitir percorrer o
fluxo de abertura até encerramento de uma solicitação e consultar seu histórico.
Os critérios e o conjunto de cenários serão acordados antes da avaliação.
O atendimento será registrado por cenário, com as falhas e limitações encontradas.

\chapter{Metodologia proposta}

\section{Fundamentação e levantamento}

O plano deverá ser apoiado por literatura e documentação técnica pertinentes
à solução escolhida. Neste protótipo, o manual é usado apenas como referência
de organização do planejamento. A base específica do tema deverá ser pesquisada,
citada no texto e acrescentada às referências antes da execução.

Prevê-se caracterizar o processo atual por documentos autorizados e, quando
cabível, consulta às pessoas envolvidas. Instrumentos, participantes, critérios
de seleção e forma de análise deverão ser definidos com o orientador. Não se
iniciará coleta que dependa de autorização ou avaliação ética antes de sua
obtenção.

\section{Desenvolvimento e avaliação}

Os requisitos serão priorizados conforme a demanda confirmada e o tempo
disponível. O desenvolvimento será organizado em incrementos: primeiro o
registro, depois o acompanhamento e, por fim, o histórico e as restrições de
acesso. Cada incremento deverá gerar uma demonstração e registro de revisão.

A avaliação prevista usará cenários funcionais e dados sintéticos. Para cada
cenário, deverão ser registrados a versão do protótipo, os passos executados,
o resultado observado e a evidência correspondente. Caso seja incorporada
avaliação com usuários, o desenho deverá ser atualizado previamente, com
instrumentos, autorizações e critérios próprios.

\section{Limites e tratamento dos resultados}

Testes em ambiente de demonstração não demonstram, isoladamente, adequação
para produção ou melhoria do serviço real. No relatório final, compare os
resultados com os objetivos e critérios acordados, expondo falhas, mudanças
de escopo e limites. Preserve a distinção entre evidência observada e benefício
esperado.

\chapter{Plano de trabalho}

\section{Etapas, responsáveis e dependências}

O estudante conduzirá levantamento, desenvolvimento e registro das evidências,
com revisão do orientador nos marcos de diagnóstico, proposta e avaliação.
O responsável pelo laboratório, caso participe, deverá autorizar o acesso
ao ambiente e às informações pertinentes. Essas participações são previstas,
sem pressupor aceite ou vínculo de pessoas reais.

O desenvolvimento depende da confirmação da demanda e da priorização dos
requisitos. A avaliação depende de versão identificada e cenários aprovados.
As semanas da \tabelaref{cronograma} são relativas ao início autorizado do
projeto e deverão ser convertidas em datas no planejamento efetivo.

\begin{undftabela}[H]{Cronograma proposto e entregas de acompanhamento}{cronograma}
  \small
  \begin{tabular}{>{\raggedright\arraybackslash}p{.18\textwidth}>{\raggedright\arraybackslash}p{.34\textwidth}>{\raggedright\arraybackslash}p{.34\textwidth}}
    \hline
    \textbf{Período} & \textbf{Etapa prevista} & \textbf{Entrega para conferência} \\
    \hline
    Semanas 1 e 2 & Fundamentação e diagnóstico autorizado & Síntese do contexto e da demanda \\
    Semanas 3 e 4 & Proposta, requisitos e desenho da avaliação & Escopo e cenários acordados \\
    Semanas 5 a 8 & Construção incremental & Versões demonstráveis do protótipo \\
    Semanas 9 e 10 & Execução dos testes e análise & Evidências, falhas e limitações \\
    Semanas 11 e 12 & Redação, revisão e entrega & Relatório final e materiais conferidos \\
    \hline
  \end{tabular}
\end{undftabela}

\section{Recursos, riscos e acompanhamento}

Prevê-se utilizar computador, ambiente de desenvolvimento e controle de versão.
Custos, licenças e acesso a infraestrutura deverão ser estimados após a escolha
tecnológica. Atraso de autorização, alteração de demanda e indisponibilidade
de ambiente são riscos a acompanhar. Se necessário, reduza o escopo com o
orientador e registre a decisão, sem substituir dados reais por coleta fictícia.

\section{Passagem ao relatório final}

Este documento termina no planejamento. Após a execução, atualize o contexto,
os requisitos e a metodologia com o percurso realizado. Acrescente o produto
entregue, os resultados da avaliação, a discussão e a parte conclusiva. Anexe
ou disponibilize materiais de apoio conforme autoria e permissões e confira
os elementos obrigatórios e opcionais da entrega definidos pelo curso.
